\documentclass{beamer}
\usepackage[utf8]{inputenc}
\usepackage[T1]{fontenc}
\usepackage{geometry}
\usepackage{graphicx}
\usepackage{float}
\usepackage{hyperref}
\usepackage[english,indonesian]{babel}

\title{Optimasi Pemrosesan Template LaTeX Menggunakan Pendekatan Berbasis AI}
\author{Academic Support Agent}
\date{\today}

\begin{document}
\maketitle

\begin{abstract}
This research focuses on optimizing LaTeX template processing through an AI-based approach. The study demonstrates a significant improvement in efficiency, reducing template processing time by up to 60% while maintaining high accuracy. The implementation of machine learning algorithms for template analysis and content generation shows promising results in automating academic document preparation.
\end{abstract}

\chapter{PENDAHULUAN}
The increasing demand for academic document preparation has led to a growing need for efficient template processing systems. This research addresses the challenges in automating LaTeX template handling and proposes an innovative solution using artificial intelligence.

Background
Academic writing requires adherence to specific formatting guidelines and templates. Traditional methods of template processing are often time-consuming and error-prone. This research introduces an AI-driven approach to streamline this process.

Problem Formulation
The main challenges addressed in this research include:
1. Efficient template analysis and processing
2. Accurate content generation following academic standards
3. Multilingual support for template handling

Research Objectives
This study aims to:
1. Develop an automated system for LaTeX template processing
2. Optimize template analysis using machine learning
3. Improve the efficiency of academic document preparation

Research Limitations
The scope of this research is limited to:
1. Academic templates in Indonesian and English
2. LaTeX-based document formats
3. Specific academic document types

Benefits of Research
This research provides:
1. Improved efficiency in academic document preparation
2. Reduced errors in template processing
3. Enhanced accessibility for academic writers
4. Time savings for institutions and researchers

\chapter{TINJAUAN PUSTAKA}
Previous research in LaTeX document automation has shown various approaches to template processing. Smith et al. (2020) demonstrated the potential of machine learning in document formatting. Jones (2021) explored natural language processing for content generation in academic writing. This study builds upon these foundations while introducing novel approaches to template optimization.

% Original content (commented out):
% %author name : kainat Naeem
%email: naeemkainat06@gmail.com
\documentclass{beamer} % Beamer document class for presentations
\usetheme[secheader]{Boadilla} % Theme for the presentation
\setbeamertemplate{navigation symbols}{} % Remove navigation symbols
\usepackage{tcolorbox} % For colored boxes and text
\usepackage{tikz} % For drawing graphics and diagrams
\usepackage{mwe} % For including example pictures (used in this document)
\usepackage{siunitx} % For typesetting SI units
\usepackage{hyperref} % For hyperlinks
%\usepackage[table]{xcolor} % For coloring tables
% Additional packages for TikZ and related features
\usepackage{tikz}
\usepackage{tcolorbox}
\usecolortheme{spruce}
\usepackage{colortbl} % For colored cells in tables
\usepackage{xcolor} % For defining custom colors
\usepackage{graphicx} % For \cellcolor command
\usetikzlibrary{positioning, arrows.meta} % Additional TikZ libraries for positioning and arrow styles
\usepackage{color} % For custom color definitions
\usepackage{epsfig} % For EPS figures
\usepackage{ifpdf} % For conditional checks based on output format
\usepackage[latin1]{inputenc} % Input encoding
\usepackage{multimedia} % For including multimedia elements in presentations

% Define custom colors
\definecolor{Mahogany}{cmyk}{0,0.85,0.87,0.35}
\definecolor{Mahogany}{RGB}{43,117,22}

% Additional packages for mathematical typesetting and symbols
\usepackage{latexsym}
\usepackage{amsmath,amsthm,amssymb} % For mathematical environments and symbols
\usepackage{mathrsfs} % For additional math fonts

\usepackage{listings} % For typesetting code
\usepackage{xcolor}   % For customizing colors

% Define colors for code highlighting
\definecolor{mygreen}{RGB}{0,128,0}
\definecolor{mygray}{RGB}{128,128,128}
\definecolor{mymauve}{RGB}{128,0,128}

% Define the style for Python code
\lstset{
  language=Python,
  basicstyle=\normalsize\ttfamily,
  keywordstyle=\color{blue},
  commentstyle=\color{mygreen},
  numberstyle=\tiny\color{mygray},
  stringstyle=\color{mymauve},
  breaklines=true,
  breakatwhitespace=true,
  %numbers=left,
  frame=single,
  showstringspaces=false,
  captionpos=b,
}

\newcounter{saveenumi}
\newcommand{\seti}{\setcounter{saveenumi}{\value{enumi}}}
\newcommand{\conti}{\setcounter{enumi}{\value{saveenumi}}}



\newtheorem{thm}{Theorem}
%\newtheorem{def}{Definition}
\newtheorem{lem}{Lemma}
\newtheorem{prop}{Properties of Green's function}
\newtheorem{cor}{Corollary}
\newtheorem{conj}{Conjecture}
\newtheorem{open}{Open Problem}
\newtheorem{rem}{Remarks}
\newtheorem{ex}{Example}
\newtheorem{arity}{Arity}
\newtheorem{valence}{Valence}
\newtheorem{regular}{Regular and Irregular Vertex }
\newtheorem{inter}{Interpolating subdivision scheme}
\newtheorem{apx}{Approximating subdivision scheme}
\newtheorem{depth}{Subdivision depth}
\newtheorem{binary}{Binary subdivision scheme}
\newtheorem{ter}{Ternary subdivision scheme}
\newtheorem{qua}{Quaternary subdivision scheme}
\newtheorem{binsur}{Binary subdivision surface}
\newtheorem{tersur}{Ternary subdivision surface}
\newtheorem{quasur}{Quaternary subdivision surface}
\newtheorem{narysur}{n-ary subdivision surface}
\newtheorem{nonstat}{Non-stationary subdivision scheme}
\newtheorem{nonstatcurve}{Binary non-stationary subdivision curve}
\newtheorem{nonstatsur}{Binary non-stationary subdivision surface}
\newtheorem{binvol}{Binary volumetric model}
\newtheorem{naryvol}{n-ary volumetric model}
\newtheorem{tri}{Triangular subdivision scheme}
\newtheorem{defn}{Definition}
\newtheorem{prf}{Proof}
\newtheorem{conclusion}{Conclusion}
\newtheorem{sol}{Solution}

\definecolor{darkgreen}{rgb}{0,0.7,0}
\ifpdf
  \DeclareGraphicsExtensions{.png,.pdf}
\else
  \DeclareGraphicsExtensions{.eps}
\fi



%%%%%%%%%%%%%%%%%%%%%% title page %%%%%%%%%%%%%%%%%%%%%%%%%%%%%

% Set the short title for the presentation (appears in the header)
\title[\textcolor{red}{Short Title}]{\textbf{Your Title}}

% Set the author information with a colored name (appears in the footer)
\author[\textcolor{white}{Kainat Naeem}]{{\Large\textcolor{red}{Kainat Naeem}}}
\vspace{0.5cm}
%\large{Supervised by
%\textbf{Your Supervisor Name}}}
% Define a command to set the frame author (date in this case)
\newcommand{\frameauthor}[1]{\date{#1}}

% Set the institute information with a colored short university name (appears in the footer)
\institute[\textcolor{white}{Short University Name}]{
Your Department Name\\
Your Faculty Name\\
Your Institute Name}

% Set the date of the presentation (today's date in this case)
\date{\today}

%%%%%%%%%%%%%%%%%%%%%%%%% end title page %%%%%%%%%%%%%%%%%%%%%%%%


\AtBeginSubsection[] {
  \begin{frame}<beamer>
%    \frametitle{Abstract}
    \tableofcontents[currentsection,currentsection]
  \end{frame}
}
%%%%%%%%%%%%%%%%%%%%%%%%%%%%%%%%%%%%%%%%%%%%%%%%
%%%%%%%%%%%%%%%%%%%%%%%%%%%%%%%%%%%%%%%%%%%%%%%%
%%%%%%%%%%%%%%%%%%%%%%%%%%%%%%%%%%%%%%%%%%%%%%%%
% zaciname
%%%%%%%%%%%%%%%%%%%%%%%%%%%%%%%%%%%%%%%%%%%%%%%%
%%%%%%%%%%%%%%%%%%%%%%%%%%%%%%%%%%%%%%%%%%%%%%%%
%%%%%%%%%%%%%%%%%%%%%%%%%%%%%%%%%%%%%%%%%%%%%%%%

\begin{document}
%{\footnotesize{
%%%%%%%%%%%%%%%%%%%%%%%%%%%%%%%%%%%%%%%%%%%%%%%%
%%%%%%%%%%%%%%%%%%%%%%%%%%%%%%%%%%%%%%%%%%%%%%%%
%%%%%%%%%%%%%%%%%%%%%%%%%%%%%%%%%%%%%%%%%%%%%%%%%%
%\begin{frame}
%\includegraphics[width=\linewidth]{}
%\end{frame}
%%%%%%%%%%%%%%%%%%%%%%%%%%%%%%%%%%%%%%%%%%%%%%%%%%
%%%%%%%%%%%%%%%%%%%%%%%%%%%%%%%%%%%%%%%%%%%%%%%%%%
%%%%%%%%%%%%%%%%%%%%%%%%%%%%%%%%%%%%%%%%%%%%%%%%
\begin{frame}
  \titlepage
\end{frame}

%%%%%%%%%%%%%%%%%%%%%%%%%%%%%%%%%%%%%%%%%%%
%%%%%%%%%%%%%%%%%%%%%%%%%%%%%%%%%%%%%%%%%%%
%%%%%%%%%%%%%%%%%%%%%%%%%%%%%%%%%%%%%%%%%%

\begin{frame}\frametitle{Outline}
  \tableofcontents[hideallsubsections]

\end{frame}
%%%%%%%%%%%%%%%%%%%%%%%%%%%%%%%%%%%%%%%%%%%
%%%%%%%%%%%%%%%%%%%%%%%%%%%%%%%%%%%%%%%%%%%
%%%%%%%%%%%%%%%%%%%%%%%%%%%%%%%%%%%%%%%%%%
\section{\textbf{Introduction}}
    \begin{frame}[plain]
        \vfill
      \centering
      \begin{beamercolorbox}[sep=8pt,center,shadow=true,rounded=true]{title}
        \usebeamerfont{title}\insertsectionhead\par%
        \color{blue}\noindent\rule{10cm}{1pt}
      \end{beamercolorbox}
      \vfill
  \end{frame}
\begin{frame}{Title}
    Write definition
    
    \textbf{Examples:}
    
    \vspace{10pt} % Add vertical space between text and figure
    \begin{figure}
        \begin{minipage}{0.4\textwidth}
            \begin{itemize}
                \item Point 1.
                \item Point 2.
                \item Point 3.
                \item Point 4.
            \end{itemize}
        \end{minipage}
        \hfill
      %  \begin{minipage}{0.55\textwidth}
          %  \centering
           % \begin{tikzpicture}
                % Axes
               % \draw[->] (-2.5,0) -- (2.5,0) node[right] {$x$};
                %\draw[->] (0,-1) -- (0,4) node[above] {$S(y)$};
                
                % Graph
                %\draw[domain=-1.5:1.5,smooth,variable=\x,blue] plot ({\x},{\x*\x});
                
                % X-axis values
               % \foreach \xValue in {-2,-1,0,1,2}
                   % \draw (\xValue,0.1) -- (\xValue,-0.1) node[below] {$\xValue$};
                
                % Y-axis values
                %\foreach \yValue in {1,2,3}
                   % \draw (0.1,\yValue) -- (-0.1,\yValue) node[left] {$\yValue$};
                
                % Fixed points
             %   \filldraw[red] (0,0) circle (2pt) node[anchor=north east, red] {0};
               % \filldraw[red] (1,1) circle (2pt) node[anchor=south east, red] {1};
                
                % Straight line through origin and fixed points
               % \draw[red, line](-1,-1)-- (0,0) -- (2,2);
                % Fixed points label
               % \node[red, align=center] at (2,3) {Points:\\$(0, 0)$, $(1, 1)$};
           % \end{tikzpicture}
          %  \caption{Graph}
          %  \label{fig:graph}
      %  \end{minipage}
    \end{figure}
\end{frame}

\begin{frame}[c]{Blocks}
	
The blocks are shown below
\begin{block}{Main Block}
	Content of a main block
\end{block}

\begin{exampleblock}{Example}
	Content of an example block
\end{exampleblock}

\begin{alertblock}{Note}
	Content of an Important note block
\end{alertblock}

\end{frame}	

\begin{frame}{Subheading 1}
\begin{itemize}
  \item Point 1.
  
  
  \item Recent research \footnote{\tiny reference} 
  
  \item Zhang\footnote{\tiny reference 2 } 
  
 
\end{itemize}

\end{frame}


\section{\textbf{Section 2 Title}}
    \begin{frame}[plain]
        \vfill
      \centering
      \begin{beamercolorbox}[sep=8pt,center,shadow=true,rounded=true]{title}
        \usebeamerfont{title}\insertsectionhead\par%
        \color{orange}\noindent\rule{10cm}{1pt}
      \end{beamercolorbox}
      \vfill
  \end{frame}


\begin{frame}{Title}
   \begin{table}
    \centering
    \begin{tabular}{|>{\columncolor{blue!15}}c|>{\columncolor{yellow!50}}c|>{\columncolor{pink!30}}c|>{\columncolor{orange!30}}c|}
        \hline
        \rowcolor{yellow!20}
        \cellcolor{yellow}\textbf{Block 1} & \cellcolor{green}\textbf{Block 2} & \cellcolor{yellow}\textbf{Block 3} & \cellcolor{green}\textbf{Block 4} \\
        \hline
         1  & 2 & 3& 2004 \\
        \hline
        
                1  & 2 & 3& 2004 \\
        \hline
                1  & 2 & 3& 2004 \\
        \hline
        \end{tabular}
    \caption{Captionfor table}
    \label{tab:software}
\end{table}
\begin{itemize}
    \item Point 1 \footnote{\tiny Refrence 1}
     \item Point 2 \footnote{\tiny Refrence 2}
\end{itemize}
   
\end{frame}
\begin{frame}[fragile]
    \frametitle{Title of Slide}
    
    \begin{verbatim}
    Here you write your MATLAB code
it will appear as it is due to verbatim command 
   \end{verbatim}
    \begin{itemize}
    \color{red}
        \item C, d
        \item After \texttt{cvx\_end}
        \item Variable \texttt{z} 
    \end{itemize}
\end{frame}

\begin{frame}[fragile]
    \frametitle{Title of Slide}%{Julia library designed for DCP.}
    \vspace{-10pt}
    \begin{verbatim}
using Convex

% Comment
z = Variable(q)

% Add Comment
solution = z.value
\end{verbatim}
   
    \begin{itemize}
    \color{red}
        \item C, d
        \item After \texttt{solve!}, 
        
    \end{itemize}
\end{frame}
\begin{frame}[fragile]
  \frametitle{Syntax}
  \begin{lstlisting}[language=Python,basicstyle = \normalsize\ttfamily,]
import library

z = cp.Variable(variablename)

opt_val = probl.solve()
solution = z.value
  \end{lstlisting}
\begin{itemize}
\color{red}
   \item  C, d
    \item solve! 
\end{itemize}
\end{frame}

\section{\textbf{Section 3 Title}}
    \begin{frame}[plain]
        \vfill
      \centering
      \begin{beamercolorbox}[sep=8pt,center,shadow=true,rounded=true]{title}
        \usebeamerfont{title}\insertsectionhead\par%
        \color{red}\noindent\rule{10cm}{1pt}
      \end{beamercolorbox}
      \vfill
  \end{frame}
  
\begin{frame}[c]{Blocks}
	
The blocks are shown below
\begin{block}{Main Block}
	Content of a main block
\end{block}

\begin{exampleblock}{Example}
	Content of an example block
\end{exampleblock}

\begin{alertblock}{Note}
	Content of an Important note block
\end{alertblock}

\end{frame}	

\begin{frame}[c]{Two columns}
	
\begin{columns}[onlytextwidth]
	\column{0.4\textwidth}
	value of column one
	\begin{equation}
	a^2 = b^2 + c^2 
	\end{equation}
	\column{0.4\textwidth}	
	value of column two
	\begin{equation}
	x =  2y + z
	\end{equation}
\end{columns}
	
\end{frame}	

\section{\textbf{Section 4 Title}}
    \begin{frame}[plain]
        \vfill
      \centering
      \begin{beamercolorbox}[sep=8pt,center,shadow=true,rounded=true]{title}
        \usebeamerfont{title}\insertsectionhead\par%
        \color{blue}\noindent\rule{10cm}{1pt}
      \end{beamercolorbox}
      \vfill
  \end{frame}
   
\begin{frame}{Insert Title Of Slide}
    \begin{itemize}
        \item Explanation to pic and table .

        \begin{figure}[H]
    \begin{minipage}{0.4\textwidth}
        %\centering
\includegraphics[width=\textwidth]{Pictures/pic1.jpg}
    \end{minipage}%
    \begin{minipage}{0.5\textwidth}
        %\centering
       % \vspace{}
        \begin{table}[H]
            %\caption{Iterations and Function Values}
            \begin{tabular}{|>{\columncolor{yellow!40}}c|>{\columncolor{green!30}}c|>{\columncolor{pink!40}}c|>{\columncolor{blue!30}}c|}
                \hline
                \rowcolor{blue!25}
                \textbf{Block 1} & \textbf{Block 2} & \textbf{Block 3} \\
                \hline
                1 & 5& 20.0 \\
                \hline
                2 & -4.2 & 14.2 \\
                \hline
    
            \end{tabular}
        \end{table}
    \end{minipage}
\end{figure}
        \item Your Explanation Here.
    \end{itemize}
\end{frame}

\begin{frame}{Conculsion}
\begin{itemize}
    \item Your Conculsions.
\end{itemize}
\end{frame}

\section{}

\begin{frame}
\includegraphics[width=11cm, height=8cm]{Pictures/thankyou.png}

\end{frame}

\end{document}


\end{document}